\chapter{绪论}

\begin{yijubox}
E1 \texttt{8-1考情分析.pptx} slide33：第一章绪论＝重点掌握基础概念。\\
E2 \texttt{9-1考点分析.pptx}：第 1--3 章偏概念，掌握基本公式与解释即可。\\
\textbf{本章结论}：不考大题，只在选择、填空、判断中出 1--3 题（约 3--9 分），
全部是"定义 + 概念辨析"。复习策略＝把定义逐字记住，把易混概念对比表背下来。
\end{yijubox}

\section{流体力学的研究任务与研究方法}

\begin{kaodian}{2}{流体力学的任务与三种研究方法}
\textbf{是什么：}流体力学研究流体的\textbf{宏观平衡和运动规律}，以及流体与固体之间的相互作用。

\textbf{教材原话（p1）：}“流体力学的任务是研究流体在平衡或运动时所遵循的基本规律及其在工程中应用的科
学，是力学的一个重要分支学科。”固体有一定的体积和形状、不易变形；液体有一定的体积而无一定的形状、
不易压缩、形状随容器形状而变、可有自由表面；气体则既无一定的体积又无一定的形状、容易压缩、
将充满整个容器、没有自由表面。液体和气体统称为流体。

\textbf{流体与固体的差别（教材 p1）：}在于对外力的抵抗能力不同。流体分子间距离较大、内聚力较小，
几乎不能承受拉力；“静止的流体则不能抵抗剪切力，即使在很小的剪切力作用下，静止流体都很容易发生
变形或流动，这种特性称为流体的易流动性。流体的易流动性是流体的基本特征。”（教材 p1）

\textbf{三种研究方法（必背，考选择/填空）：}教材的提法是\textbf{理论分析、科学实验、数值模拟}三种
（教材 p1）。注意填空要按教材原文写\textbf{“科学实验”“数值模拟”}，不要写成“实验研究”“数值计算”。
\begin{center}
\begin{tabularx}{\textwidth}{@{}lXX@{}}
\toprule
方法 & 做法 & 优缺点 \\
\midrule
理论分析 & 抓住主要因素建立模型 $\to$ 建立封闭方程组与边界/初始条件 $\to$ 解析求解 & 推导严谨、答案精确；但只适用于较简单的理论模型，复杂流动无能为力 \\
科学实验（实验研究） & 按相似原理设计模型，实测相似准则数，拟合准则方程式推广 & 更接近实际、结果可靠；但受测量精度与模型符合程度限制，有些流动无法实验 \\
数值模拟（数值计算） & 确定数学模型 $\to$ 选计算方法 $\to$ 编程上机 $\to$ 分析结果 & 可解解析法不能解的问题，省时省钱；但依赖理论分析与实验提供模型，且难涵盖全部影响因素 \\
\bottomrule
\end{tabularx}
\end{center}
三种方法相辅相成：理论指导实验与数值计算，实验检验理论、提供建模依据，数值计算弥补前两者的不足。

\textbf{常考题型：}填空（"流体力学的研究方法有\rule{2cm}{0.4pt}、\rule{2cm}{0.4pt}、\rule{2cm}{0.4pt}三种"）；选择（问哪一种方法"答案精确"或"更接近实际"）。

\textbf{重要度 \xstarfull{2}}\quad\yiju{E1 slide33 只要求"基础概念"；E4 历年真题中未见独立成题}\quad\nandu{易}

\begin{banfa}
三方法优缺点用\textbf{六个字}记：理论"准而窄"，实验"真而难"，数值"广而依"。
选择题问"哪种方法对复杂流动无能为力"$\to$理论分析；问"最接近实际"$\to$实验研究。
\end{banfa}
\end{kaodian}

\section{连续介质模型}

\begin{kaodian}{4}{连续介质假设}
\textbf{是什么：}把流体看作由\textbf{连续分布的质点}组成的介质，其物理量（密度、速度、压强等）
是空间坐标和时间的\textbf{连续函数}。

\textbf{为什么引入：}真实流体由离散分子组成、分子间有空隙；但工程问题的特征尺度远大于分子平均自由程，
因此可以"抹平"分子间隙，用连续函数和微积分来描述流体运动——这是整个流体力学数学工具的地基。

\textbf{适用条件：}所研究问题的特征尺度远大于分子平均自由程（一般工程问题均满足）；
高空稀薄气体、真空技术等不适用。

\textbf{教材原话（p2）：}“流体的连续介质模型假定流体是由连续分布的流体质点所组成，即认为流体所占据的
空间完全由没有任何空隙的流体质点所充满，流体质点在时间过程中做连续运动。”
教材把“流体质点”定义为“流体中宏观尺寸非常小而微观尺寸又足够大的任意一个物理实体”，有三个特点：
① 宏观尺寸非常小、无尺度、可视为一个点，而微观尺寸足够大、内含足够多的流体分子；
② 具有质量、密度、压强、流速、动能等宏观物理量，这些物理量“是流体质点中大量流体分子的统计平均值”；
③ 流体质点的形状可任意划定。

教材给出引进该模型的依据：流体分子平均自由程“往往远小于一般工程问题的特征尺寸”，
且流体力学关心的是“流体宏观特性，即大量分子的统计平均特性”（教材 p2）。
于是“表征流体性质和运动特性的物理量和力学量一般为空间坐标和时间变量的连续函数，
如压强 $p=p(x,y,z,t)$”（教材 p2），这样才能用数学分析方法研究流体运动。

\begin{zhenti}{2022 年 818 第 1 题}{单项选择题}
连续介质假设意味着（\quad）。\\
(A) 流体分子互相紧连\quad (B) 流体的物理量是连续函数\quad (C) 流体分子间有孔隙\quad (D) 流体不可压缩
\end{zhenti}
\noindent\textbf{答案：(B)。}连续介质假设的本质是"物理量连续"，不是"分子紧挨着"。

\textbf{常考题型：}选择题（如上）、填空题（"连续介质假设认为流体是由\underline{\hspace{1.5cm}}组成的介质"）。
\textbf{重要度 \xstarfull{4}}\quad\yiju{E4 2022 年 818 第 1 题直接考；E1 slide33"第一章＝重点掌握基础概念"}\quad\nandu{易}

\begin{yicuo}
把"连续介质"理解成"分子紧挨着、没有间隙"是\textbf{错}的。假设针对的是\emph{宏观物理量}的连续性，
它并不否认分子的离散性；"(C) 流体分子间有孔隙"恰恰是事实，所以不能选。
\end{yicuo}
\end{kaodian}

\section{流体的主要物理性质}

\begin{kaodian}{3}{密度、比容与重度}
\begin{gongshi}{密度的定义（教材式 1.1、1.2）}
\[ \rho=\frac{m}{V}\qquad\qquad \rho_A=\lim_{\Delta V\to0}\frac{\Delta m}{\Delta V} \]
\noindent 依次为\textbf{教材式 1.1}（均质流体密度）与\textbf{教材式 1.2}（非均质流体中点 A 的密度，教材 p3）。\\
\textbf{含义：}单位体积的流体所具有的质量称为流体的密度（教材 p2 原话）。对非均质流体，
若包含 A 点的微元体积 $\Delta V$ 中的流体质量为 $\Delta m$，
则“流体中 A 点的密度”按教材式 1.2 取 $\Delta V\to0$ 的极限，得到 A 点的\textbf{点密度}。\\
\textbf{适用条件：}连续介质假设成立。教材 p3 原话：“一般来说，流体的密度随压强和温度而变化。
不过，液体的密度随压强和温度的变化很小，通常情况下可视为常数”，常用值：水 $\rho=1000\ \mathrm{kg/m^3}$，
水银 $\rho=13600\ \mathrm{kg/m^3}$（教材 p3）。
\end{gongshi}

\begin{gongshi}{比容与重度（教材式 1.3、1.4）}
\[ u=\frac{1}{\rho}\ \ (\mathrm{m^3/kg}),\qquad\qquad \gamma=\frac{G}{V}=\rho g\ \ (\mathrm{N/m^3}) \]
\noindent 依次为\textbf{教材式 1.3}、\textbf{教材式 1.4}（教材 p3）。\\
\textbf{含义：}比容＝单位质量流体所占体积（密度的倒数）；重度＝单位体积流体的重量。
教材 p3 原话：“密度的倒数，即单位质量的流体所具有的体积称为比容”“单位体积的流体所具有的重量
称为流体的重度”，教材式 1.4 中 $G$ 为流体重量、$V$ 为流体体积。\\
\textbf{适用条件：}$g$ 取当地重力加速度，工程计算常取 $g=9.8\ \mathrm{m/s^2}$（个别题取 9.81）；
重度与密度用教材式 1.4 互换时单位必须统一为 $\mathrm{N/m^3}$ 与 $\mathrm{kg/m^3}$。
\end{gongshi}

\textbf{常考题型：}填空题小计算——"已知密度 $\rho=851\ \mathrm{kg/m^3}$，则重度 $\gamma=\underline{\quad}$、比容 $u=\underline{\quad}$"。
\textbf{重要度 \xstarfull{3}}\quad\yiju{E4 2015 年试题填空第 1 题（由密度与运动黏度求重度与动力黏度）；E5 题库同类题}\quad\nandu{易}

\begin{banfa}
$\gamma=\rho g$ 一定要带 $g$，写成 $\gamma=\rho$ 就是单位（N/m³ vs kg/m³）没检查。
填空题给 $\nu$ 求 $\mu$ 时用 $\mu=\nu\rho$；给 $\rho$ 求 $\nu$ 时用 $\nu=\mu/\rho$。
\textbf{三步自检：}① 单位换算成国际单位；② 查表值是否与题给温度一致；③ 结果数量级是否合理（20 ℃ 水的 $\nu\approx1.0\times10^{-6}\ \mathrm{m^2/s}$）。
\end{banfa}
\end{kaodian}

\begin{kaodian}{3}{流体的压缩性与膨胀性}
\begin{gongshi}{体积压缩系数、体积弹性模量、体积膨胀系数（教材式 1.5--1.8）}
\[ k=-\frac{1}{V}\frac{\dd V}{\dd p}\ \ (\mathrm{Pa^{-1}}),\qquad\qquad
   k=\frac{1}{\rho}\frac{\dd\rho}{\dd p}\ \ (\mathrm{Pa^{-1}}) \]
\[ K=\frac{1}{k}=\rho\frac{\dd p}{\dd\rho}\ \ (\mathrm{Pa}),\qquad\qquad
   \alpha_V=\frac{1}{V}\frac{\dd V}{\dd T}\ \ (\mathrm{K^{-1}}) \]
\noindent 依次为\textbf{教材式 1.5}、\textbf{教材式 1.6}、\textbf{教材式 1.7}、\textbf{教材式 1.8}（教材 p3--4）。
教材式 1.5 与 1.6 都是体积压缩系数，前者用体积变化表示、后者用密度变化表示：
由 $\dd m=\dd(\rho V)=\rho\,\dd V+V\dd\rho=0$ 得 $\dfrac{\dd\rho}{\rho}=-\dfrac{\dd V}{V}$，代入即得教材式 1.6。\\
\textbf{含义：}$k$＝单位压强增量引起的体积相对减小量（右侧加负号是为了让它为正）；
$K$＝$k$ 的倒数，$K$ 越大压缩性越小；$\alpha_V$＝单位温升引起的体积相对增大量。\\
\textbf{适用条件：}$k$、$\alpha_V$ 都是“定温/定压”下定义的；三者都随流体种类、温度、压强变化。
教材 p3 原话：“任何流体都是可压缩的，只是可压缩程度有所不同而已”“当流体的压缩性对所研究的流动
影响不大时，忽略其压缩性，这样的流体称为不可压缩流体，不可压缩流体是理想化的力学模型”。
\end{gongshi}

\textbf{结论（背）：}任何流体都是可压缩的，只是程度不同。
\textbf{液体}在一般平衡与运动问题中压缩性可忽略、按不可压缩处理；
但\textbf{水击}和\textbf{水中爆炸}问题中必须按可压缩流体处理。
\textbf{低速气体}也可视为不可压缩流体。
教材 p4 原话（膨胀性）：“通常液体的体积膨胀系数很小，一般工程问题中当温度变化不大时，
可不予考虑，而气体的体积膨胀系数却很大。”

\textbf{常考题型：}选择/判断（"下列哪种情况必须考虑液体的压缩性"$\to$水击）；
填空（"体积弹性模量越大，流体压缩性越\underline{\ \ 小\ }"）。
\textbf{重要度 \xstarfull{3}}\quad\yiju{E5 题库选择题（工程概念题）；E6 教材 1.3.2 节为独立小节}\quad\nandu{易}

\begin{yicuo}
易错：把 $k$ 说成"压缩系数越大越难压缩"。恰恰相反——$k$ 越大越\textbf{容易}压缩；
$K=1/k$ 越大才越\textbf{难}压缩。另一处：1.5 式的负号是"约定"，不是物理上体积会变大。
\end{yicuo}
\end{kaodian}

\begin{kaodian}{5}{黏性——牛顿内摩擦定律、动力黏度与运动黏度}
\textbf{是什么：}流体运动时，质点之间发生相对运动就会产生抵抗相对运动的\textbf{切向阻力}，
这种性质叫黏性。黏性是流体产生阻力和能量损失的根源。
教材 p4 原话：“静止的流体没有抵抗剪切变形的能力，而运动的流体，当流体质点之间发生相对运动时，
质点之间就会产生切向阻力（摩擦阻力）抵抗其相对运动……这种特性称为流体的黏性”；
教材还称黏性为“流体的重要属性”，指出它“是流体运动中产生阻力和能量损失的原因”。

\begin{gongshi}{牛顿内摩擦定律（教材式 1.9、1.10）}
\[ T=\mu A\frac{\dd u}{\dd y},\qquad\qquad \tau=\frac{T}{A}=\mu\frac{\dd u}{\dd y} \]
\noindent 依次为\textbf{教材式 1.9}、\textbf{教材式 1.10}（教材 p4--5）。教材 p5 原话：
“式（1.9）和式（1.10）称为牛顿内摩擦定律”。\\
\textbf{含义：}教材 p4 原话：“实验证明，流体内摩擦阻力 $T$ 的大小与速度法向梯度 $\dfrac{\dd u}{\dd y}$
和接触面积 $A$ 成正比，并与流体的性质有关”，比例系数就是动力黏度 $\mu$；
单位面积上的内摩擦力即为\textbf{切应力} $\tau$。\\
\textbf{适用条件：}只适用于\textbf{牛顿流体}的\textbf{层流}（平行层状）运动；$\dfrac{\dd u}{\dd y}$ 必须按速度剖面的法向取。
\textbf{线形速度分布（教材 p5）：}“当两平板间距离 $h$ 和速度 $U$ 不大时，速度 $u$ 沿其法线方向呈线性分布”，
此时 $u=\dfrac{U y}{h}$、$T=\mu A\dfrac{U}{h}$、$\tau=\mu\dfrac{U}{h}$（$y$ 自静止板算起）。
\end{gongshi}

\begin{gongshi}{动力黏度与运动黏度（教材式 1.11）}
\[ \mu\ \ (\mathrm{Pa\cdot s}),\qquad\qquad \nu=\frac{\mu}{\rho}\ \ (\mathrm{m^2/s}) \]
\noindent 运动黏度的定义式为\textbf{教材式 1.11}（教材 p5。教材 p5 原话：
“在流体力学中，动力黏度 $\mu$ 经常与流体密度 $\rho$ 结合在一起以 $\mu/\rho$ 的形式出现。
将这个比值定义为运动黏度 $\nu$”）。\\
\textbf{含义：}$\mu$ 反映流体内摩擦的强弱，直接出现在牛顿内摩擦定律里；
$\nu=\mu/\rho$ 只是因为它在雷诺数等无量纲数中经常作为一个整体出现，才单独定义，\textbf{并不表示黏性大小}。\\
\textbf{适用条件：}温度、压强一定。工程常用范围内黏度受压强影响小，\textbf{主要随温度变化}；
教材给出 $\nu$ 的单位为 $\mathrm{m^2/s}$。
\end{gongshi}

\textbf{速度梯度的物理意义：}取矩形微元面，上下层速度差 $\dd u$，经 $\dd t$ 发生角变形，
角变形速度 $\dfrac{\dd\theta}{\dd t}=\dfrac{\dd u}{\dd y}$——所以\textbf{速度梯度＝角变形速度}。

\textbf{温度对黏度的影响（必考）：}
\begin{center}
\begin{tabularx}{\textwidth}{@{}lXX@{}}
\toprule
流体 & 黏性来源 & 温度升高时 \\
\midrule
液体 & 分子间\textbf{内聚力}（分子间距小，内聚力强）为主 & 分子间距加大 $\to$ 内聚力减小 $\to$ 黏度\textbf{减小} \\
气体 & 分子热运动的\textbf{动量交换}（内聚力可忽略）为主 & 热运动加剧 $\to$ 黏度\textbf{增大} \\
\bottomrule
\end{tabularx}
\end{center}

\textbf{教材 p6 原话（黏性机理）：}“液体分子间距小，内聚力强，黏性作用主要来源于分子内聚力，
当液体温度升高时，其分子间距加大，内聚力减小，黏度随温度上升而减小；而气体和液体不同，
气体的内聚力极小，可以忽略，其黏性作用可以说完全是分子热运动中动量交换的结果”，
所以气体黏度随温度升高而\textbf{增大}，与液体相反。

\textbf{教材表 1.1（标准大气压、不同温度时水和空气的黏度）：}
\begin{center}
\begin{tabular}{@{}rcccc@{}}
\toprule
$t$\,/\,℃ & $\mu_{\text{水}}$ & $\nu_{\text{水}}$ & $\mu_{\text{空气}}$ & $\nu_{\text{空气}}$ \\
 & $10^{-3}\ \mathrm{Pa\cdot s}$ & $10^{-6}\ \mathrm{m^2/s}$ & $10^{-5}\ \mathrm{Pa\cdot s}$ & $10^{-6}\ \mathrm{m^2/s}$ \\
\midrule
0 & 1.792 & 1.792 & 1.72 & 13.7 \\
10 & 1.308 & 1.308 & 1.78 & 14.7 \\
20 & 1.005 & 1.007 & 1.83 & 15.7 \\
30 & 0.801 & 0.804 & 1.87 & 16.6 \\
40 & 0.656 & 0.661 & 1.92 & 17.6 \\
50 & 0.549 & 0.556 & 1.96 & 18.6 \\
60 & 0.469 & 0.477 & 2.01 & 19.6 \\
70 & 0.406 & 0.415 & 2.04 & 20.6 \\
80 & 0.357 & 0.367 & 2.10 & 21.7 \\
90 & 0.317 & 0.328 & 2.16 & 22.9 \\
100 & 0.284 & 0.296 & 2.18 & 23.6 \\
\bottomrule
\end{tabular}
\end{center}
\textbf{读表要点：}水（液体）的 $\mu$ 随温度升高而减小、空气（气体）的 $\mu$ 随温度升高而增大，
与上表“水减气增”一致；20 ℃ 水的 $\nu\approx1.007\times10^{-6}\ \mathrm{m^2/s}$，可作填空题查表值。

\begin{tishi}
\textbf{教材表 1.1 的一处疑问（两说并列）：}该表水 $\mu$ 列的表头在教材原文中作 $10^{-2}\ \mathrm{Pa\cdot s}$，
而表内数值（0 ℃ 为 1.792、20 ℃ 为 1.005 等）属 $10^{-3}\ \mathrm{Pa\cdot s}$ 量级；若真按 $10^{-2}$ 读，
水的黏度会比公认值大 10 倍。故本表按 $10^{-3}\ \mathrm{Pa\cdot s}$ 列出，同时保留教材表头一说；
空气 $\mu$ 列表头为 $10^{-3}\ \mathrm{Pa\cdot s}$、表值 0.0172 等，本表已换算到 $10^{-5}\ \mathrm{Pa\cdot s}$ 的同一量级。
答题时按数值量级判断即可。
\end{tishi}

\textbf{教材同型题（教材第 1 章无编号例题）：}教材 p8 习题 1.5——平板尺寸 $0.8\ \mathrm{m}\times0.2\ \mathrm{m}$，
在油面上以 $U=1\ \mathrm{m/s}$ 运动，油膜厚 $\delta=1\ \mathrm{mm}$、$\mu=1.15\ \mathrm{Pa\cdot s}$，求黏性力，
解法即本节 $\tau=\mu\dfrac{U}{\delta}$、$T=\tau A$，与上面 banfa ③ 三步走完全同型。

\begin{zhenti}{2026 年真题回忆版}{简答题}
概念题中出现"理想流体"；计算题第 1 题为\textbf{牛顿流体}（斜板上的静压强/切应力问题）。
\end{zhenti}

\textbf{常考题型：}选择（"汽油属于（牛顿/非牛顿）流体"、黏度随温度的变化方向）；
填空（"运动黏度的量纲是\rule{2cm}{0.4pt}"$\to$ $\mathrm{L^2/T}$）；
简答（黏性与流动性的关系）；计算（斜板间流体求 $\tau$ 或求 $\mu$）。
\textbf{重要度 \xstarfull{5}}\quad\yiju{E4 2026 计算题第 1 题（牛顿流体）；E5 题库选择第 2 题"汽油＝牛顿流体"、问答题第 1 题"流动性与黏滞性"；E4 常以填空考 $\mu$、$\nu$ 计算；E1 slide16 填空考"名词定义"}\quad\nandu{易}

\begin{banfa}
\textbf{①温度方向"水减气增"记法}：水（液体）越热越"稀"，空气越热越"稠"。\\
\textbf{②$\nu$ 不表示黏性大小}：选项里出现"运动黏度越大黏性越大"一定是错的。
$\nu=\mu/\rho$ 还随 $\rho$ 变，物理上没有"黏性"的绝对含义。\\
\textbf{③斜板/平板类计算三步走}：(a) 求两板间距 $h$；(b) 速度线性分布时 $\dfrac{\dd u}{\dd y}=\dfrac{U}{h}$；
(c) 代 $\tau=\mu U/h$ 或 $T=\mu A U/h$。注意题目给的是 \emph{动力}黏度还是 \emph{运动}黏度，给 $\nu$ 要先乘 $\rho$。\\
\textbf{④量纲速判}：$\mu\to\mathrm{M/(L\cdot T)}$（即 $\mathrm{Pa\cdot s}$），$\nu\to\mathrm{L^2/T}$，$\tau\to\mathrm{Pa}$。
\end{banfa}
\end{kaodian}

\begin{kaodian}{3}{牛顿流体、理想流体与非牛顿流体}
\textbf{牛顿流体：}切应力与角变形速度（速度法向梯度）成\textbf{正比}，即遵守牛顿内摩擦定律的流体。
水、空气、汽油等属牛顿流体。教材 p6 原话：“凡作用在流体上的切应力与它所引起的角变形速度
（速度法向梯度）成正比，即遵守牛顿内摩擦定律的流体称为牛顿流体；否则，称为非牛顿流体”，
并明确“本书只讨论牛顿流体，而非牛顿流体是流变学的研究对象”。

\textbf{非牛顿流体：}不遵守牛顿内摩擦定律，$\tau\sim\dfrac{\dd u}{\dd y}$ 关系不是过原点的直线。
\begin{center}
\begin{tabularx}{\textwidth}{@{}lXl@{}}
\toprule
类型（教材图 1.2 中的线） & 特征 & 举例 \\
\midrule
理想宾汉流体（B 线） & 切应力需达到某一值才开始剪切变形，之后变形速度为常数 & 泥浆、血浆 \\
伪塑性流体（C 线） & 黏度随角变形速度增大而\textbf{减小} & 尼龙、颜料、油漆 \\
膨胀性流体（D 线） & 黏度随角变形速度增大而\textbf{增大} & 生面团、浓淀粉糊 \\
\bottomrule
\end{tabularx}
\end{center}

\textbf{理想流体：}不具黏性的流体，是现实中不存在的\textbf{假想流体}。
教材 p6 原话：“实际流体都具有黏性。不具有黏性的流体称为理想流体，它是客观世界中并不存在的一种假想的流体。”
引入目的是简化分析：当实际流体黏度小或所研究区域速度梯度小、黏滞力远小于其他力时可以忽略黏性；
即使对黏性占主导的问题，也可先研究理想流体、再研究实际流体。

\textbf{常考题型：}选择（判断某流体属哪一类；"理想流体是指（\quad）"）；
判断（"理想流体是黏度很小的流体"$\to$ 错，是\emph{不具黏性}的假想流体）。
\textbf{重要度 \xstarfull{4}}\quad\yiju{E5 题库选择第 2 题（汽油＝牛顿流体）；E4 2026 回忆版简答涉及"理想流体"；E6 教材 1.3.3 节}\quad\nandu{易}

\begin{yicuo}
"理想流体＝黏度很小的流体"错；"理想流体＝不可压缩流体"也错。
理想流体的\textbf{唯一}含义是"无黏性"；不可压缩是另一条独立假设，两者不能混为一谈。
\end{yicuo}
\end{kaodian}

\begin{kaodian}{2}{液体的表面张力}
\textbf{是什么：}液体自由表面上相邻部分之间的\textbf{拉力}，方向与液面相切、与两相邻部分的分界线垂直。
教材 p6 原话：“液体的表面张力是液体自由表面上相邻部分之间的拉力，其方向与液面相切，并与两相邻部分的分界线垂直。”

\textbf{本质：}液面上分子受到液体内部分子的吸引与上方气体分子的吸引不平衡，合力垂直液面指向液体内部，
使液体有尽量缩小表面的趋势，宏观上表现为表面拉力。

\begin{gongshi}{表面张力系数}
\[ \sigma=\frac{F}{L}\ \ (\mathrm{N/m}) \]
\textbf{含义：}作用在单位长度上的表面张力。$\sigma$ 与液体性质、纯度、温度及接触介质有关。
教材 p6 原话：“表面张力一般产生在液体和气体相接触的自由表面上，也可以产生于液体与固体的接触面上
或与另一种液体的接触面上。”\\
\textbf{适用条件：}教材 p6 又指出“表面张力仅在液体的自由表面存在，液体内部并不存在，所以它是一种
局部受力现象”；因其很小，一般对液体宏观运动不起作用，可以忽略。
\end{gongshi}

\begin{tishi}
教材 p6 前后两处表述并不完全一致：一处说表面张力“也可以产生于液体与固体的接触面上或与另一种液体的
接触面上”，另一处说“仅在液体的自由表面存在”。两说并列记忆，应试以“\textbf{液体内部不存在、
只存在于界面（自由表面或相界面）上}”为准。
\end{tishi}

\textbf{教材表 1.2（几种液体与空气接触时的表面张力系数）：}
\begin{center}
\begin{tabular}{@{}lcc@{}}
\toprule
液体 & $t$\,/\,℃ & $\sigma$\,/\,$(\mathrm{N/m})$ \\
\midrule
水 & 20 & 0.07275 \\
水银 & 20 & 0.465 \\
酒精 & 20 & 0.0223 \\
四氯化碳 & 20 & 0.0257 \\
丙酮 & 16.8 & 0.02344 \\
甘油 & 20 & 0.065 \\
苯 & 20 & 0.0289 \\
润滑油 & 20 & $0.025\sim0.035$ \\
\bottomrule
\end{tabular}
\end{center}
\textbf{读表要点：}水银的表面张力系数约为水的 $6.4$ 倍，这也是水银在毛细管中不润湿玻璃、液面明显下降的原因。

\textbf{常考题型：}选择/判断（表面张力存在于液体内部？$\to$ 错）；
填空（表面张力系数的单位是 N/m）。
\textbf{重要度 \xstarfull{2}}\quad\yiju{E4 历年真题未独立成题；E6 教材 1.3.4 节}\quad\nandu{易}
\end{kaodian}

\begin{kaodian}{2}{毛细现象}
\textbf{是什么：}把细管插入液体后，管内液面高于或低于管外液面的现象（教材 p7）。
教材 p7 用管内液体的\textbf{内聚力}和它与管壁之间的\textbf{附着力}的相对大小来解释。

\textbf{判据（教材 p7）：}附着力大于内聚力时液体润湿管壁，接触角 $\theta$ 为锐角
（水与玻璃约为 $8.5^\circ$），管内液面升高并呈凹形；附着力小于内聚力时液体不润湿管壁，
接触角 $\theta$ 为钝角（水银与玻璃约为 $140^\circ$），管内液面下降并呈凸形。

\begin{gongshi}{毛细现象中液柱上升（或下降）的高度（教材式 1.12）}
\[ h=\frac{4\sigma\cos\theta}{\rho g d}\ \ (\mathrm{m}) \]
\noindent 该式为\textbf{教材式 1.12}（教材 p7）。\\
\textbf{含义：}$\sigma$ 为液体的表面张力系数，$\theta$ 为接触角，$\rho$ 为液体密度，$d$ 为细管直径。
教材 p7 原话：“液柱上升或下降的高度与管径成反比，$d$ 小，$h$ 则大。”\\
\textbf{适用条件：}细管内的液体与管壁润湿（或完全不润湿）且管径足够小；接触角须按具体液—固组合取值，
$\theta$ 小于 $90^\circ$ 对应液面升高、大于 $90^\circ$ 对应液面下降（式中 $\cos\theta$ 自带符号）。
\end{gongshi}

\textbf{工程意义（教材 p7）：}由教材式 1.12 可知管径越小液位差越大，故使用液位计、单管测压计时
应选择合适管径，以免毛细现象造成读数误差。
\textbf{重要度 \xstarfull{2}}\quad\yiju{E6 教材 1.3.4（2）毛细现象（教材式 1.12）；E4 历年真题未独立成题}\quad\nandu{易}
\end{kaodian}

\section{作用在流体上的力}

\begin{kaodian}{3}{质量力与表面力}
\textbf{是什么：}作用在流体上的力按作用方式分为两类。教材 p8 原话：
“质量力指作用于流体块中各流体质点的非接触性外力，例如重力、惯性力等。质量力与流体块质量成正比，
又称为体积力。作用于单位质量流体上的质量力叫单位质量力。”“表面力指作用于流体块表面上的外力。
表面力可按其作用方向分为垂直并指向流体块表面的压力以及与表面平行的切向力。”

\begin{gongshi}{质量力与表面力}
\[ \bF_{\text{质量}} = m\bm{a}\ \ \text{（如重力、惯性力、离心力）},\qquad
   \bF_{\text{表面}} = \text{压力} + \text{切应力（摩擦力）} \]
\textbf{含义：}\textbf{质量力}（又称体积力）作用于流体块内\textbf{每一个质点}，大小与质量成正比，
是\textbf{非接触力}；\textbf{表面力}作用于流体与流体（或流体与固体）的\textbf{接触面}上，
大小与面积成正比，可分解为法向的\textbf{压力}与切向的\textbf{切应力}。\\
\textbf{适用条件：}质量力是\emph{外}力且为非接触力；表面力必须通过接触面传递。
\end{gongshi}

\begin{gongshi}{质量力的分量式与单位质量力（教材式 1.13、1.14）}
\[ \bF=F_x\bm{i}+F_y\bm{j}+F_z\bm{k},\qquad\qquad
   \bm{f}=\frac{\bF}{m}=f_x\bm{i}+f_y\bm{j}+f_z\bm{k}\ \ (\mathrm{m/s^2}) \]
\noindent 依次为\textbf{教材式 1.13}、\textbf{教材式 1.14}（教材 p8）。\\
\textbf{含义：}质量力用三个分量表示（教材式 1.13）；教材式 1.14 是作用于单位质量流体上的质量力，
其量纲与加速度相同，$f_x,f_y,f_z$ 一般是坐标与时间的函数。\\
\textbf{适用条件：}分量式对直角坐标系成立；质量力有势时可用势函数 $W$ 表示三个分量。
\end{gongshi}

\begin{gongshi}{表面力的压强与切应力（教材式 1.15、1.16）}
\[ p=\lim_{\Delta A\to0}\frac{\Delta P}{\Delta A}\ \ (\mathrm{Pa}),\qquad\qquad
   \tau=\lim_{\Delta A\to0}\frac{\Delta T}{\Delta A}\ \ (\mathrm{Pa}) \]
\noindent 依次为\textbf{教材式 1.15}、\textbf{教材式 1.16}（教材 p8）。\\
\textbf{含义：}$\Delta P$、$\Delta T$ 分别是微元面积 $\Delta A$ 上的法向压力与切向力；
$p$ 是单位面积上的法向压力（压强），$\tau$ 是单位面积上的切向力（切应力）。\\
\textbf{适用条件：}必须取 $\Delta A\to0$ 的极限才是“一点”的应力；工程计算中常近似取有限面积上的平均值。
\end{gongshi}

\begin{jielun}
\textbf{教材 p8 结论：}“理想流体或者静止流体的表面应力只有压强而无切应力。”
这正是第 2 章静压强两个重要特性（方向垂直指向作用面、同一点各方向压强相等）的来源。
\end{jielun}

\textbf{单位质量力：}把质量力除以质量，得单位质量力 $f=(f_x,f_y,f_z)$，其量纲与加速度相同。
\textbf{有势的质量力：}若存在势函数 $W$ 使 $f_x=-\dfrac{\partial W}{\partial x}$（等号可整体变号），
则称该质量力有势——重力就是有势力，$W=gz$。

\textbf{常考题型：}选择（"体积力有势"的含义；下列哪种力是质量力）；
填空（质量力的量纲、单位质量力的表达式）。
\textbf{重要度 \xstarfull{3}}\quad\yiju{E5 题库选择第 3 题（体积力有势）；E6 教材 1.4 节；质量力是流体平衡微分方程的直接输入（第 2 章）}\quad\nandu{易}

\begin{banfa}
判别口诀：\textbf{"接触的是表面力，不接触的是质量力"}。
题目里"作用于流体块内每个质点上"$\to$质量力；"作用于接触面上"$\to$表面力。
另外记：压力是表面力的\textbf{法向}分量，切应力是\textbf{切向}分量——
第 2 章静力学中流体不能承受切应力，所以静止流体表面力只剩压力，这正是静压强特性的来源。
\end{banfa}
\end{kaodian}

\section{本章小结}

\begin{center}
\begin{tabularx}{\textwidth}{@{}lccX@{}}
\toprule
考点 & 重要度 & 难度 & 一句话抓住 \\
\midrule
研究任务与三种方法 & \xstarfull{2} & 易 & 理论分析 / 科学实验 / 数值模拟（准-真-广） \\
连续介质假设 & \xstarfull{4} & 易 & 物理量是连续函数（$\neq$分子紧挨着） \\
密度、比容、重度 & \xstarfull{3} & 易 & $\gamma=\rho g$，$\mu=\nu\rho$ \\
压缩性与膨胀性 & \xstarfull{3} & 易 & $K=1/k$；水击必须考虑液体压缩性 \\
黏性与牛顿内摩擦定律 & \xstarfull{5} & 易 & $\tau=\mu\,\dd u/\dd y$；水减气增 \\
牛顿/理想/非牛顿流体 & \xstarfull{4} & 易 & 理想流体＝\emph{无黏性}的假想流体 \\
表面张力 & \xstarfull{2} & 易 & 界面存在，液体内部没有（教材式：$\sigma=F/L$） \\
毛细现象 & \xstarfull{2} & 易 & 教材式 1.12：$h=4\sigma\cos\theta/(\rho gd)$，$h\propto 1/d$ \\
质量力与表面力 & \xstarfull{3} & 易 & 接触＝表面力，不接触＝质量力 \\
\bottomrule
\end{tabularx}
\end{center}
